\section{Erd\H{o}s Problem 88: every small induced edge count}

\subsection{1) Statement and external results}
For each fixed $\epsilon>0$, does there exist $\delta>0$ such that
every $n$-vertex graph with no clique or independent set of size
at least $\epsilon\log n$ has an induced subgraph with exactly
$m$ edges for every integer $0\leq m\leq\delta n^2$?
This is the Erd\H{o}s--McKay conjecture, proved by Kwan, Sah,
Sauermann and Sawhney. We reconstruct its deduction from their
stronger interval theorem and the standard density theorem they
state, with all uniformity issues explicit.

Call a graph $C$-Ramsey when all its cliques and independent sets
have size at most $C\log_2 n$. The external inputs are:
\begin{enumerate}
\item For each $C>0$ there is $\rho(C)>0$ such that all sufficiently
large $C$-Ramsey graphs satisfy
$e(G)\geq\rho(C)\binom n2$. This is the Erd\H{o}s--Szemer\'edi
density theorem, stated as Theorem 4.1 in the cited paper.
\item For each $C>0$ and $\eta>0$, every sufficiently large
$C$-Ramsey graph has an induced subgraph with exactly $m$ edges
for every integer $0\leq m\leq(1-\eta)e(G)$.
This is Theorem 1.1 of Kwan--Sah--Sauermann--Sawhney.
\end{enumerate}
Their proofs, including the necessary local probability estimates,
are external dependencies rather than claims proved in this note.

\subsection{2) Deduction, including the small orders}
Take $\log$ in the problem to mean natural logarithm and choose
$C=\epsilon\log2$. The hypothesis implies the $C$-Ramsey
condition. Let $n_0\geq2$ be a common threshold for the two inputs,
the second with $\eta=1/2$, and put $\rho=\rho(C)$.
For $n\geq n_0$,
\[
 \frac12e(G)\geq\frac{\rho}{2}\binom n2
               \geq\frac{\rho n^2}{8}.
\]
Choose
\[
 \delta=\min\left(\frac\rho8,\frac1{2n_0^2}\right)>0.
\]
Every integer $m\leq\delta n^2$ then lies in the interval covered
by the second theorem. For $1\leq n<n_0$, we have
$\delta n^2<1$, so the only nonnegative integer to consider is
$m=0$, realized by a one-vertex induced subgraph. Thus the same
$\delta$, depending only on $\epsilon$, works at every order for
which the hypothesis is met. If zero-vertex graphs are permitted,
their empty induced graph treats that case as well.

\subsection{3) Why the exact interval input is essential}
For comparison, deleting vertices from any graph gives a decreasing
sequence of induced edge counts from $e(G)$ to $0$, with each jump
at most $n-1$. Hence for any target $m$ there is an induced subgraph
whose edge count is within $n-1$ of $m$: take the first step crossing
the target. This elementary argument does not ensure equality.
Nor does knowing that there are many distinct induced edge counts
ensure every integer in an initial interval is present. The stronger
external theorem addresses precisely that missing exactness.

\textbf{Outcome: SOLVED, established affirmative resolution
reconstructed relative to the two named external theorems.}
No new proof of their analytic ingredients or local formal
verification is claimed.

\subsection{4) Sources}
\begin{itemize}
\item M. Kwan, A. Sah, L. Sauermann and M. Sawhney,
\emph{Anticoncentration in Ramsey graphs and a proof of the
Erd\H{o}s--McKay conjecture}, Theorems 1.1 and 4.1,
\texttt{https://arxiv.org/abs/2208.02874}.
\item T. F. Bloom, Problem 88,
\texttt{https://www.erdosproblems.com/88}, checked 23 September 2026.
\end{itemize}
The estimate measures coverage of the problem using these explicit
external results.
\subsection{5) COMPLETION ESTIMATE}
\textbf{COMPLETION: 100\%}
